% Ganti perintah \lipsum dengan isi intisari.
Abstrak ditulis dengan maksimum 200 kata yang disusun dalam tiga alinea. Alinea pertama berisi tentang latar belakang secara ringkas. Penggunaan \textit{floating nuclear power plant} (FNPP) atau reaktor nuklir terapung berpotensi digunakan untuk pemenuhan listrik dengan mobilitas tinggi. Reaktor KLT-40S memiliki spesifikasi data termal 150 MWt. Penelitian ini akan melakukan analisis terhadap performansi neutronik teras reaktor. 

Alinea kedua berisi tentang metodologi secara ringkas. Model desain KLT-40S dalam kode simulasi \textit{Serpent} dibuat untuk pengamatan nilai \textit{burnup} dan panjang siklus operasi teras reaktor. Variasi untuk analisis \textit{burnup} dan panjang siklus dilakukan pada beragam pengayaan dan keberadaan konten silumin pada pin bahan bakar dan peracun dapat bakar. Dilakukan analisis \textit{inherent safe} pada beberapa parameter termohidrolik. Pada analisis \textit{burnup} diamati pula dinamika pembentukan inventaris isotop \textsuperscript{239}Pu sebagai profil hambatan proliferasi. 

Alinea ketiga berisi tentang hasil dan kesimpulan secara ringkas. Panjang siklus yang mendekati data teknis dicapai pada pengayaan 15,7 wt\% dengan matriks silumin dengan panjang siklus 846,37 hari. Pada penggunaan pengayaan 18,6 wt\% dengan matriks silumin diperoleh koefisien reaktivitas suhu bahan bakar -2,68092 pcm/K, suhu pendingin sebesar -21,63755 pcm/K, tanpa penggunaan silumin diperoleh koefisien reaktivitas fraksi void -0,2154/K. Fraksi massa isotop \textsuperscript{239}Pu pada pengayaan 18,6 wt\% dengan matriks silumin diperoleh pada nilai 54,219\%.

\KataKunci{kata kunci 1, kata kunci 2, kata kunci 3, kata kunci 4}
